
    \hfill % 两张图片之间的间距
   问题三的核心目标是综合考虑BMI、孕周、年龄、身高、体重等多因素影响，建立多目标优化模型，为男胎孕妇的NIPT检测提供最优分组策略和检测时点建议。整体解决思路是基于问题二构建的广义可加混合模型，充分估计各因素对Y染色体浓度的影响，并综合估计结果的影响，给出以BMI为分组依据的检测策略。
\begin{figure}[htpb]
    \centering
    \includegraphics[width=0.8\textwidth]{figures/5.3/mindmap.png}
    \caption{问题三求解流程}
    \label{fig:}
\end{figure}
\subsection{数据可视化与作图分析}

图\ref{fig:img11123}的每个小图展示两个变量的散点关系，直观显示 BMI、孕周、体重、身高、年龄、GC 含量等变量之间的相关性与分布形态。部分变量关系近似线性（如 BMI 与体重），部分存在非线性与离群点，揭示了多因素间的潜在的相互关系。整体上，在BMI为25-40之间，样本分布趋于集中，特别是在BMI为30-35区间，体重主要分布在0.5至1.0的范围内，反映了研究样本中较高BMI群体的特征。进一步分析显示，Y染色体浓度与BMI呈负相关关系，较高的BMI通常伴随较低的Y染色体浓度，这为NIPT的时点选择提供了重要数据支持。

通过小提琴图\ref{fig:sub3},\ref{fig:sub4}的分析，可以看出数据中的体重、孕妇BMI、13号、18号和21号染色体的Z值呈现出不同程度的分布特点。BMI的数据分布相对较为集中，而染色体的Z值则显示出不同样本间的基因波动，尤其在13号染色体Z值上最为明显。结合年龄、身高等变量的分布情况，BMI和染色体Z值之间的潜在关系值得进一步分析，以确定影响NIPT准确性的关键因素。

\subsubsection{身高体重的正交化}

由于 BMI 与身高（HT）、体重（WT）存在强相关性（\(BMI\approx WT/HT^2\)），若直接纳入模型会导致共线性，干扰 BMI 的核心效应。因此需要分别计算两指标对BMI的残差，然后剥离BMI对两指标的解释部分（正交化处理）：
\begin{equation}
\begin{aligned}
    rHT_i &= \log(HT_i) - \hat{\mathbb{E}}[\log(HT) \mid \log(BMI)]_i.\\
    rWT_i &= \log(WT_i) - \hat{\mathbb{E}}[\log(WT) \mid \log(BMI), rHT]_i.\\
\end{aligned}
\end{equation}

使rHT和rWT保留了独立于 BMI 的身材信息，既补充了变量维度，又避免干扰 BMI 的核心作用。


\begin{figure}[htbp]
    \centering
    % 第一行
    \begin{minipage}{0.45\textwidth}
        \centering
        \includegraphics[width=\linewidth]{figures/5.3/01.png}
        \captionof{figure}{各指标间散点图}
        \label{fig:img11123}
    \end{minipage}
    \hfill
    \begin{minipage}{0.45\textwidth}
        \centering
        \includegraphics[width=\linewidth]{figures/5.3/02.png}
        \captionof{figure}{各指标间热力图}
        \label{fig:sub2}
    \end{minipage}

    % 换行
    \vspace{1em} % 可选，添加一些垂直间距

    % 第二行
    \begin{minipage}{0.45\textwidth}
        \centering
        \includegraphics[width=\linewidth]{figures/5.3/03.png}
        \captionof{figure}{13、18、21号染色体小提琴图}
        \label{fig:sub3}
    \end{minipage}
    \hfill
    \begin{minipage}{0.45\textwidth}
        \centering
        \includegraphics[width=\linewidth]{figures/5.3/04.png}
        \captionof{figure}{孕妇各指标小提琴图}
        \label{fig:sub4}
    \end{minipage}

    \label{fig:main}
\end{figure}


\subsubsection{GAMM模型构建}
基于问题一所构建的原来的 GAMM 模型，本部分上新增年龄和相对 BMI 正交变化后的身高与体重变量构建GAMM模型，其固定效应包括模型的群体平均趋势 \(\mu_i(t)\) ，由多个平滑函数 \(f_k(\cdot)\) 组合而成，形式为 \(\mu_i(t)=\underbrace{f_1(\text{GA}_c) \times f_2(\text{BMI}_c)}_{\text{主交互}}+f_3(r\text{HT}_c)+f_4(r\text{WT}_c)+f_5(\text{Age}_c)\)，其中各 \( f_k(\cdot) \) 为三次回归样条（Age 取低自由度以抑制过拟合），“\( \cdot_c \)”代表中心化。
随机截距：
\[
\text{logit}(Y_i^*(t)) = \mu_i(t) + b_i + \varepsilon, \quad b_i \sim \mathcal{N}(0, \sigma_b^2), \, \varepsilon \sim \mathcal{N}(0, \sigma_\varepsilon^2).
\]


\noindent $\bullet$ \textbf{达标概率与累计曲线}
设定达标阈值 \(\theta = 0.04\)，个体在时刻 t 的一次达标概率 \(p_i(t)=\mathbb{P}(Y_i(t)\geq\theta) \approx \sigma(\mu_i(t)-\ell_\theta)\)；累计达标曲线为\(F_i(t)=\max_{s\leq t}p_i(s)\)表示截至时刻 t，个体能达到的最大达标概率。 为了让结果更稳健，引入总方差 \(\sigma^2=\sigma_b^2+\sigma_\varepsilon^2\)，重新计算一次达标概率 \(p_i^{(\text{err})}(t)=\mathbb{P}(\mu_i(t)+Z\geq\ell_\theta)=1 - \Phi\left(\frac{\ell_\theta - \mu_i(t)}{\sigma}\right)\)，最后再基于 \(p_i^{(\text{err})}(t)\) 计算累计达标曲线 \(F_i^{(\text{err})}(t)=\max_{s\leq t}p_i^{(\text{err})}(s)\)。

\subsubsection{动态规划最优分组算法}

基于 GAMM 模型输出，计算个体条件达标概率 \( p_i(t) = \text{logit}^{-1}(\mu_i(t) - \eta^{\dagger}) \) 及其累计形式 \( F_i(t) = \max_{s \leq t} p_i(t) \)。定义风险函数 \( \mathcal{R}_i(t) = \sum_{k=0}^{\infty} w(t_k) \left[ F_i(t_k) - F_i(t_{k-1}) \right] + \lambda \left[ 1 - F_i(t) \right] \)，其中权重参数取 \( (W_1, W_2, W_3) = (0, 2, 10) \)，首检失败成本 \( \lambda = 0.5 \)。采用动态规划算法进行最优分割，以 \( \text{Crit}(k) = \text{DP}_k[N] + \gamma (k - 1) \log N \) 为选择准则（\( \gamma = 0.15 \)），确保分组结果既最小化总体风险又控制模型复杂度。

\subsubsection{求解结果}  

GAMM 残差方差 \(0.0625\)，随机截距方差 \(0.1554\)。结构惩罚下得到最优组数 \(G^{*} = 3\)，组间 BMI 边界与指标如下（表\ref{tab:divide}）。组 1 与组 3 的最佳时点为 11.75 周，组 2 为 23.75 周；对应一次达标率分别为 0.613、0.710、0.589（重测率为 0.387、0.290、0.411）。

\begin{table}[htbp]
  \centering
  \caption{分组检测相关指标}
  \label{tab:divide}
  \begin{tabular}{ccccccc}
    \toprule
    组别 & 人数 & BMI范围 & 最优检测时点(周) & 一次达标率 & 重测率  \\
    \midrule
    1 & 136 & 20.70 \~{}31.34& 11.75 & 0.6124 & 0.3876  \\
    2 & 51 & 31.40 \~{} 33.29 & 23.75 & 0.7103 & 0.2897  \\
    3 & 80 & 33.05 \~{} 46.87 & 11.75 & 0.5886 & 0.4114  \\
    \bottomrule
  \end{tabular}
\end{table}
